\documentclass[../DoAn.tex]{subfiles}
\begin{document}

% =========================================================
\section{Mở đầu chương}
% =========================================================

Chương này trình bày các nền tảng lý thuyết phục vụ trực tiếp cho việc xây
dựng hệ thống trợ lý học vụ dựa trên kiến trúc \textit{Retrieval-Augmented
Generation} (RAG) --- tức kiến trúc sinh câu trả lời có bổ trợ truy hồi. Khác
với các bài toán hỏi đáp thông thường, dữ liệu học vụ có tính đặc thù cao:
mỗi quy định thường gắn với phạm vi áp dụng cụ thể, thời điểm hiệu lực, mã
ngành và mã học phần nhất định. Do đó, mô hình ngôn ngữ lớn (LLM --- Large
Language Model) không thể chỉ dựa vào tri thức tham số học được trong quá
trình huấn luyện --- hệ thống cần có cơ chế truy hồi tài liệu có kiểm soát
trước khi đưa ra câu trả lời.

Về bố cục, nội dung chương được tổ chức theo trình tự từ bài toán đến giải
pháp kỹ thuật. Phần đầu phân tích đặc thù của dữ liệu và truy vấn học vụ.
Tiếp theo, chương trình bày nguyên lý RAG, quy trình xử lý tài liệu và phân
đoạn nội dung (chunking), biểu diễn ngữ nghĩa bằng embedding vector, tìm kiếm
từ khóa với BM25, truy hồi lai kết hợp Qdrant và Elasticsearch, lọc theo siêu
dữ liệu, xếp hạng lại bằng mô hình cross-encoder, hiểu truy vấn và định tuyến,
mô hình RAG có tác nhân (agentic RAG) theo kiến trúc Planner--Executor, cùng
các cơ chế hỗ trợ như tìm kiếm web dự phòng và đánh giá hệ thống. Toàn bộ
các khái niệm được trình bày ở mức tổng quát, đặt nền tảng lý thuyết cho phần
đề xuất giải pháp và triển khai ở các chương sau.

% =========================================================
\section{Bối cảnh bài toán hỏi đáp học vụ}
% =========================================================

\subsection{Đặc thù nguồn tri thức học vụ}

Nguồn tri thức học vụ không phải là một tập văn bản đồng nhất. Thực tế, nó
bao gồm nhiều nhóm tài liệu có cấu trúc và mục đích khác nhau: chương trình
đào tạo, quy định học vụ, kế hoạch học kỳ, biểu mẫu thủ tục và các dịch vụ
hỗ trợ sinh viên. Vì mỗi nhóm có tiêu chí truy hồi riêng, hệ thống phân chia
dữ liệu thành bốn kho tri thức theo miền --- bao gồm chương trình đào tạo, quy
định học vụ, kế hoạch và thông báo, cùng sổ tay hỗ trợ sinh viên --- như đã
định nghĩa trong Bảng~\ref{tab:ch1-collections} ở Chương~1.

Việc phân chia theo miền như vậy giúp giảm nhiễu đáng kể trong quá trình truy
hồi. Nếu toàn bộ tài liệu được gộp vào một chỉ mục chung, những đoạn có từ
khóa tương tự nhưng khác mục đích sẽ dễ cạnh tranh nhầm với nhau trong danh
sách kết quả --- chẳng hạn, một thông báo về kế hoạch học kỳ có thể lấn át
thông tin quy định tốt nghiệp khi sinh viên hỏi về điều kiện ra trường.

\subsection{Đặc điểm truy vấn của sinh viên}

Truy vấn học vụ thường ngắn, mang tính tự nhiên và không phải lúc nào cũng
trùng khớp với thuật ngữ trong văn bản gốc. Trong thực tế, sinh viên hay hỏi
theo cách diễn đạt đời thường, ví dụ ``em cần điều kiện gì để tốt nghiệp'',
``môn này học kỳ mấy'' hay ``ngành Việt Nhật có khác IT1 không''. Bên cạnh đó,
truy vấn còn có thể chứa các thực thể cụ thể như mã ngành, mã học phần, khóa
sinh viên hoặc học kỳ nhất định. Đặc biệt trong các cuộc hội thoại nhiều lượt,
câu hỏi sau thường dùng đại từ và tham chiếu ngầm kiểu ``ngành của em'', ``môn
đó'' hay ``so với chương trình kia''.

Những đặc điểm này cho thấy bài toán hỏi đáp học vụ không đơn thuần là tìm
kiếm văn bản. Hệ thống phải phối hợp đồng thời retrieval theo ngữ nghĩa,
retrieval chính xác theo từ khóa, lọc theo siêu dữ liệu và xử lý ngữ cảnh hội
thoại. Tuy nhiên, lịch sử hội thoại và thông tin hồ sơ cá nhân chỉ nên được
sử dụng khi câu hỏi hiện tại thực sự phụ thuộc vào chúng. Nếu không kiểm soát
tốt, hệ thống có thể rơi vào hiện tượng context bleed (*ngữ cảnh bị nhiễm
chéo*) --- tức là mang sai thông tin ngành hoặc khóa từ lượt hỏi trước sang
câu hỏi mới, dẫn đến kết quả truy hồi bị lệch hướng.

\subsection{Yêu cầu đối với câu trả lời có căn cứ}

Trong miền học vụ, một câu trả lời thoạt nghe hợp lý nhưng thiếu căn cứ có
thể khiến sinh viên đưa ra quyết định sai, chẳng hạn đăng ký nhầm môn hoặc
bỏ lỡ thời hạn quan trọng. Do đó, hệ thống hỏi đáp cần đáp ứng một số yêu
cầu cốt lõi: truy hồi tài liệu liên quan trước khi sinh câu trả lời; ưu tiên
văn bản còn hiệu lực và tránh dùng văn bản đã bị thay thế; xử lý được các
tham chiếu điều khoản kiểu ``Điều~5, Khoản~1''; thừa nhận khi ngữ cảnh không
đủ thông tin; và hỗ trợ kiểm tra hồi quy để phát hiện lỗi ở từng tầng xử lý.

% =========================================================
\section{Nguyên lý Retrieval-Augmented Generation}
% =========================================================

\subsection{Giới hạn của tìm kiếm từ khóa và mô hình sinh thuần túy}

Tìm kiếm từ khóa truyền thống như TF--IDF hay BM25 có thế mạnh với các chuỗi
chính xác --- mã học phần, số điều khoản, tên biểu mẫu --- nhưng dễ bỏ sót tài
liệu khi người dùng diễn đạt khác so với văn bản gốc. Ở chiều ngược lại, LLM
có khả năng hiểu ngôn ngữ tự nhiên và tổng hợp câu trả lời mạch lạc, nhưng lại
không đảm bảo nắm được các văn bản nội bộ mới nhất và dễ sinh nội dung không
được kiểm chứng --- hiện tượng thường gọi là hallucination (*sinh nội dung
không có căn cứ*).

RAG giải quyết hạn chế của cả hai hướng bằng cách tách quá trình trả lời thành
hai bước rõ ràng: đầu tiên truy hồi các đoạn tài liệu phù hợp từ kho tri thức
đã được lập chỉ mục, sau đó đưa các đoạn đó vào prompt để mô hình sinh câu trả
lời bám theo nội dung đã cho. Nhờ kiến trúc tách biệt này, tri thức của hệ
thống có thể được cập nhật đơn giản bằng cách cập nhật chỉ mục mà không cần
huấn luyện lại mô hình.

\subsection{Vòng đời của một hệ thống RAG}

Nhìn tổng thể, một hệ thống RAG hoàn chỉnh gồm ba pha chính:

\begin{enumerate}
    \item \textbf{Lập chỉ mục (Indexing)}: chuyển đổi tài liệu thô, làm sạch
    nội dung, phân đoạn thành các chunk nhỏ, tạo embedding vector rồi lưu văn
    bản cùng siêu dữ liệu vào các vector database và elasticsearch.
    \item \textbf{Retrieval}: hiểu truy vấn của người dùng, chọn kho tri thức
    phù hợp, áp dụng bộ lọc siêu dữ liệu, lấy các ứng viên từ tìm kiếm vector
    và tìm kiếm từ khóa, hợp nhất kết quả rồi xếp hạng lại để chọn ra nguồn
    tốt nhất.
    \item \textbf{Sinh câu trả lời (Generation)}: định dạng ngữ cảnh đã truy
    hồi, sinh câu trả lời có căn cứ, đồng thời trả về thông tin nguồn để người
    dùng kiểm chứng.
\end{enumerate}

Điểm cần nhấn mạnh là chất lượng RAG phụ thuộc vào toàn bộ chuỗi xử lý, không
chỉ riêng mô hình sinh. Nói cách khác, nếu chunk bị phân chia sai, siêu dữ
liệu thiếu hoặc bước reranking loại nhầm tài liệu, thì mô hình có thể tạo ra
câu trả lời sai dù bản thân nó có năng lực tốt.

\subsection{Grounded generation}

Grounded generation là nguyên tắc yêu cầu mô hình chỉ trả lời dựa trên ngữ
cảnh được cung cấp, không tự suy diễn vượt quá nguồn và phải thông báo khi
không có đủ thông tin. Trong hệ thống này, tầng sinh câu trả lời cung cấp giao
diện chung cho cả chế độ thông thường và phát trực tuyến (streaming); mô hình
ngôn ngữ sử dụng có thể là DeepSeek, Gemini hoặc các điểm cuối tương thích
chuẩn OpenAI. Nhờ lớp trừu tượng này, toàn bộ pipeline không phụ thuộc vào
một thư viện cụ thể nào, trong khi vẫn đảm bảo nguyên tắc cốt lõi: bước sinh
câu trả lời luôn là bước cuối cùng, sau khi retrieval và xử lý hậu kỳ đã hoàn
tất.

% =========================================================
\section{Xử lý tài liệu và phân đoạn nội dung}
% =========================================================

\subsection{Chuyển đổi và làm sạch tài liệu}

Tài liệu đầu vào của hệ thống có thể ở nhiều dạng: PDF, DOCX, Markdown hoặc
JSON từ công cụ thu thập dữ liệu tự động (crawler). Bước chuyển đổi sang định
dạng Markdown cần đảm bảo giữ được các cấu trúc quan trọng như tiêu đề phân
cấp, danh sách, bảng và số thứ tự điều khoản --- vì nếu những cấu trúc này bị
mất, các bước chunking và retrieval về sau sẽ khó xác định đúng ngữ cảnh cho
từng đoạn. Sau khi chuyển đổi, bước làm sạch loại bỏ mục lục, dòng trang trí
và khoảng trắng dư thừa, đồng thời cố gắng giữ lại bảng ở dạng Markdown chuẩn.

\subsection{Chiến lược phân đoạn theo loại tài liệu}

Chunk là đơn vị nội dung cơ bản được lưu trữ, mã hóa thành embedding vector và
truy hồi. Chunk quá ngắn thường mất ngữ cảnh, còn chunk quá dài lại chứa nhiều
ý không liên quan, vừa làm giảm độ chính xác khi reranking vừa tiêu tốn ngân
sách context window của mô hình. Một chunk phù hợp cần bao trọn một đơn vị
nghĩa tương đối độc lập và giữ đủ siêu dữ liệu để truy vết ngược về nguồn gốc.

Do mỗi loại tài liệu học vụ có đặc điểm riêng, hệ thống áp dụng nhiều chiến
lược chunking khác nhau thay vì tách cố định theo số ký tự. Cụ thể, văn bản
quy định cần giữ cấu trúc Chương--Điều--Khoản; chương trình đào tạo cần bảo
toàn nguyên vẹn bảng học phần; thông báo cần giữ ngày tháng và tiêu đề; sổ tay
sinh viên cần giữ các bước thủ tục theo đúng trình tự.
Bảng~\ref{tab:chunking_strategies} tóm tắt các chiến lược chính.

\begin{table}[H]
\centering
\begin{tabular}{|l|p{4.0cm}|p{6.0cm}|}
\hline
\textbf{Chiến lược} & \textbf{Dữ liệu phù hợp} & \textbf{Ý tưởng chính} \\
\hline
Phân đoạn đệ quy & Chương trình đào tạo và văn bản Markdown tổng quát &
Tách theo tiêu đề, đoạn, dòng; bảo vệ bảng học phần; tạo quan hệ cha--con. \\
\hline
Phân đoạn pháp quy & Văn bản quy định (đã qua nhận dạng ký tự) & Nhận diện
Chương, Điều, Khoản và Phụ lục bằng quy tắc dựa trên chuỗi mẫu; tạo chunk
tiêu đề, chunk cha và chunk con; có xử lý dự phòng khi thiếu cấu trúc pháp lý. \\
\hline
Phân đoạn kế hoạch & Kế hoạch và thông báo thu thập tự động & Tách theo bài,
mục đánh số, kèm ngày đăng và siêu dữ liệu nguồn. \\
\hline
Phân đoạn sổ tay & Sổ tay và tài liệu hỗ trợ sinh viên & Tách theo mục, tiểu
mục và các bước thủ tục. \\
\hline
\end{tabular}
\caption{Các chiến lược chunking nội dung trong hệ thống}
\label{tab:chunking_strategies}
\end{table}

\subsection{Quan hệ cha--con và siêu dữ liệu}

Văn bản học vụ thường có cấu trúc phân cấp, do đó cách chia chunk cũng cần
phản ánh điều này. Cụ thể, chunk cha giữ ngữ cảnh rộng hơn như tiêu đề hoặc
toàn bộ nội dung một Điều, trong khi chunk con chứa nội dung trả lời trực tiếp
câu hỏi. Chính sách lập chỉ mục phân biệt rõ hai mục đích: các chunk ở dạng
tiêu đề được loại khỏi danh sách kết quả chính để tránh chiếm vị trí trong
top-$k$, nhưng vẫn được lưu trong kho để sau bước reranking có thể truy xuất
bổ sung ngữ cảnh thông qua ID chunk cha khi cần.

Ngoài nội dung văn bản, mỗi chunk còn mang theo siêu dữ liệu cấu trúc như mã
ngành, khóa áp dụng, loại tài liệu, ngày tháng, nguồn và ID tài liệu cha. Đặc
biệt, đối với tài liệu quy định, thông tin về quan hệ kế thừa --- tức mối quan
hệ thay thế giữa các phiên bản văn bản --- được lưu riêng để hệ thống có thể
nhận biết và tránh sử dụng văn bản cũ khi đã có phiên bản mới hơn.

% =========================================================
\section{Biểu diễn ngữ nghĩa bằng vector}
% =========================================================

\subsection{Embedding vector và độ đo cosine}

Embedding (*biểu diễn vector*) là phép ánh xạ chuyển văn bản thành một vector
số trong không gian nhiều chiều. Khi hai đoạn văn có ý nghĩa gần nhau, các
vector tương ứng cũng thường nằm gần nhau theo độ đo cosine:

\begin{equation}
\operatorname{cos}(u,v) =
\frac{u \cdot v}{\|u\| \|v\|}
\end{equation}

Trong giai đoạn lập chỉ mục, mỗi chunk văn bản được chuyển thành embedding
vector rồi lưu vào kho vector. Đến lúc xử lý truy vấn, câu hỏi của người dùng
cũng được mã hóa theo cùng cách, sau đó so khớp với các vector tài liệu để tìm
ra những ứng viên có ý nghĩa gần nhất. Tìm kiếm vector đặc biệt hữu ích khi
người dùng diễn đạt khác hoàn toàn so với văn bản gốc, nhưng lại không đủ mạnh
với các trường hợp cần khớp chính xác theo ký tự --- đây chính là lý do hệ
thống cần kết hợp thêm tìm kiếm từ khóa.

\subsection{Mô hình mã hóa BGE-M3 và E5 đa ngôn ngữ}

Hệ thống sử dụng đồng thời hai mô hình mã hóa đa ngôn ngữ để tăng độ bao phủ
khi truy hồi. Mô hình BGE-M3 \cite{bgem3} tạo embedding 1024 chiều và còn được dùng làm đặc
trưng cho một số thành phần phân loại miền. Mô hình E5 \cite{e5model} đa ngôn ngữ cũng tạo
embedding 1024 chiều, sử dụng tiền tố ``query:'' cho truy vấn và ``passage:''
cho tài liệu theo đặc tả riêng của mô hình này.

Hai embedding này được lưu riêng biệt trong Qdrant dưới dạng hai vector có tên
thay vì trộn sẵn ngay từ khâu lập chỉ mục. Khi truy vấn, điểm tìm kiếm vector
là tổ hợp có trọng số của hai mô hình (mặc định cân bằng $0{,}5$--$0{,}5$); nhờ
lưu tách biệt, hệ thống có thể điều chỉnh mức đóng góp của từng mô hình mà không
phải lập chỉ mục lại.

\subsection{Kho vector Qdrant}

Qdrant được sử dụng để lưu trữ embedding vector và siêu dữ liệu theo từng kho
tri thức. Mỗi điểm dữ liệu tương ứng với một chunk văn bản, bao gồm hai vector
1024 chiều và sử dụng khoảng cách cosine làm độ đo. Trường payload (*dữ liệu
mở rộng*) chứa nội dung văn bản, siêu dữ liệu, thông tin nguồn và ID chunk cha.
Ngoài khả năng tìm kiếm vector, Qdrant còn hỗ trợ lọc theo trường payload để
phối hợp với tìm kiếm từ khóa từ Elasticsearch trong luồng truy hồi lai.

% =========================================================
\section{Tìm kiếm từ khóa và BM25}
% =========================================================

\subsection{Vai trò của sparse retrieval}

Sparse retrieval (*tìm kiếm thưa, dựa trên từ vựng*) giữ lại tín hiệu từ vựng
chính xác mà dense retrieval (*tìm kiếm dày, dựa trên vector*) có thể bỏ qua.
Trong dữ liệu học vụ, các ký hiệu như ``IT-E6'', ``MI1111'', ``K70'' hay
``Điều~5'' mang ý nghĩa rất cụ thể. Nếu hệ thống chỉ dùng tìm kiếm ngữ nghĩa,
các ký hiệu này có thể bị làm mờ trong không gian vector khi cạnh tranh với
những đoạn có ý nghĩa gần nhưng không đúng thực thể mà sinh viên đang hỏi. Vì
vậy, Elasticsearch được tích hợp để thực hiện BM25, tìm kiếm theo siêu dữ liệu
và tra cứu theo định danh.

\subsection{Công thức BM25}

BM25 \cite{bm25} là phương pháp tính điểm liên quan giữa tài liệu $D$ và truy vấn $Q =
\{q_1, q_2, ..., q_n\}$ theo công thức:

\begin{equation}
\operatorname{BM25}(D,Q) =
\sum_{i=1}^{n}
\operatorname{IDF}(q_i)
\frac{f(q_i,D)(k_1+1)}
{f(q_i,D)+k_1\left(1-b+b\frac{|D|}{avgdl}\right)}
\end{equation}

Trong đó, $f(q_i,D)$ là tần suất xuất hiện của từ $q_i$ trong tài liệu, $|D|$
là độ dài tài liệu, $avgdl$ là độ dài trung bình toàn corpus, $k_1$ điều khiển
độ bão hòa tần suất và $b$ điều khiển mức chuẩn hóa theo độ dài tài liệu.
Thành phần IDF (nghịch đảo tần suất tài liệu) được tính như sau:

\begin{equation}
\operatorname{IDF}(q_i) =
\ln\left(
\frac{N-n(q_i)+0.5}{n(q_i)+0.5}+1
\right)
\end{equation}

Elasticsearch trong hệ thống được cấu hình với bộ phân tích tiếng Việt hỗ trợ
tách từ phù hợp, cùng các bước xử lý chuyển về chữ thường, đồng nghĩa và chuẩn
hóa ký tự ASCII. Bên cạnh tìm kiếm toàn văn trên nội dung, các trường như mã
ngành, mã học phần và ngày tháng được lưu dạng từ khóa để phục vụ lọc chính xác.

% =========================================================
\section{Truy hồi lai và hợp nhất kết quả}
% =========================================================

\subsection{Lý do kết hợp dense retrieval và sparse retrieval}

Dense retrieval và sparse retrieval là hai phương pháp bổ sung cho nhau. Dense
retrieval tìm được các đoạn cùng nghĩa dù cách diễn đạt có khác nhau, trong
khi sparse retrieval giữ được kết quả khớp chính xác cho mã ngành, mã môn, số
điều khoản và các cụm từ đặc thù. Chính sự kết hợp này tạo ra khả năng bao phủ
toàn diện hơn so với từng phương pháp riêng lẻ.

Trong thực tế triển khai, hệ thống chạy tìm kiếm vector và tìm kiếm từ khóa
song song cho từng kho tri thức, sau đó hợp nhất kết quả trên toàn bộ các kho
đã được chọn.

\subsection{Hợp nhất theo thứ hạng đảo nghịch (RRF)}

Do điểm cosine và điểm BM25 nằm trên các thang đo khác nhau, việc cộng trực
tiếp hai loại điểm là không phù hợp. Phương pháp hợp nhất mặc định của hệ thống
là Hợp Nhất Theo Thứ Hạng Đảo Nghịch (Reciprocal Rank Fusion --- RRF \cite{rrf2009}): thay vì
dựa vào độ lớn điểm gốc, RRF kết hợp kết quả dựa trên thứ hạng của mỗi tài liệu
trong từng tập riêng lẻ. Dạng cơ bản của RRF được tính như sau:

\begin{equation}
\operatorname{RRF}(d) =
\sum_{r \in R}
\frac{1}{k + \operatorname{rank}_r(d)}
\end{equation}

Trong đó $\operatorname{rank}_r(d)$ là thứ hạng của tài liệu $d$ trong tập kết
quả $r$, và $k$ là hằng số làm mượt. Trên cơ sở đó, hệ thống dùng một biến thể
có trọng số nhằm vẫn ưu tiên tín hiệu ngữ nghĩa mà không đánh mất tính bền vững
của RRF:

\begin{equation}
S(d) =
w_v \cdot \frac{1}{k + \operatorname{rank}_v(d)}
+ w_k \cdot \frac{1}{k + \operatorname{rank}_k(d)}
+ B_{\text{recency}}(d)
\end{equation}

Ở đây $\operatorname{rank}_v$ và $\operatorname{rank}_k$ lần lượt là thứ hạng
của tài liệu trong danh sách vector và danh sách từ khóa, $w_v$ và $w_k$ là
trọng số tương ứng (cấu hình mặc định $w_v = 0{,}8$, $w_k = 0{,}2$ và hằng số
$k = 10$), còn $B_{\text{recency}}$ là phần cộng thêm dành cho tài liệu thuộc
kho kế hoạch có ngày đăng mới hơn. Cách hợp nhất theo thứ hạng đặc biệt hữu ích
khi điểm của hai phương pháp khó so sánh trực tiếp, đồng thời miễn nhiễm với
trường hợp một tập kết quả lấn át tập kia chỉ vì chênh lệch thang điểm. Ngoài
ra, khi truy vấn chứa mã học phần hoặc yêu cầu tra bảng chính xác, tỷ lệ từ khóa
được tự động tăng lên để ưu tiên các kết quả khớp đúng định danh.

\subsection{Phương án hợp nhất tuyến tính theo điểm chuẩn hóa}

Bên cạnh RRF, hệ thống còn cài đặt một phương án hợp nhất tuyến tính dựa trên
độ lớn điểm sau chuẩn hóa, được dùng để đối chứng trong thực nghiệm. Ở phương
án này, mỗi điểm trong từng tập kết quả được chia cho giá trị lớn nhất của tập
đó (chuẩn hóa theo giá trị cao nhất, max-normalization); so với chuẩn hóa
min-max thông thường, cách này giúp tài liệu có điểm thấp nhất vẫn giữ được một
phần đóng góp khác không. Điểm hợp nhất khi đó là tổ hợp tuyến tính:

\begin{equation}
S(d) =
w_v \cdot \widehat{S_v}(d)
+ w_k \cdot \widehat{S_k}(d)
+ B_{\text{recency}}(d)
\end{equation}

trong đó $\widehat{S_v}$ và $\widehat{S_k}$ là điểm vector và điểm từ khóa sau
chuẩn hóa. Điểm khác biệt là phương án tuyến tính phản ánh trực tiếp khoảng cách
điểm giữa các tài liệu, trong khi RRF chỉ quan tâm đến thứ hạng. Kết quả so sánh
thực nghiệm giữa hai chiến lược hợp nhất này được trình bày chi tiết trong
Chương~6.

% =========================================================
\section{Truy hồi nhận biết siêu dữ liệu}
% =========================================================

\subsection{Lọc theo siêu dữ liệu và kiểm tra hiệu lực}

Lọc theo siêu dữ liệu giúp thu hẹp không gian tìm kiếm và tăng xác suất lấy
đúng văn bản mà sinh viên cần. Luồng xử lý cơ bản diễn ra theo các bước: tạo
bộ lọc theo kho tri thức, dùng Elasticsearch tìm định danh phù hợp dựa trên
siêu dữ liệu, chuyển danh sách định danh sang điều kiện lọc cho Qdrant, rồi
chạy tìm kiếm vector và từ khóa trên tập đã được lọc. Nếu bộ lọc quá chặt
khiến kết quả rỗng, hệ thống có cơ chế tự động nới lỏng điều kiện để tránh
trả lời ``không có thông tin'' quá sớm.

Sau giai đoạn truy hồi, hệ thống thực hiện thêm hai bước xử lý hậu kỳ quan
trọng. Thứ nhất, ValidityFilter (*bộ lọc hiệu lực*) đọc thông tin quan hệ kế
thừa giữa các văn bản để loại bỏ những chunk thuộc văn bản đã bị thay thế, với
điều kiện vẫn còn đủ nguồn sau khi lọc. Thứ hai, ReferenceResolver (*bộ giải
tham chiếu*) phát hiện các dẫn chiếu như ``Điều~5'' hoặc ``Khoản~1 Điều~5'',
rồi tìm chunk tương ứng bằng tra cứu siêu dữ liệu; nếu không tìm được, hệ
thống chuyển sang tìm kiếm ngữ nghĩa có ràng buộc cùng nguồn.

\subsection{Mở rộng ngữ cảnh từ chunk cha và HyDE dự phòng}

Mở rộng ngữ cảnh từ chunk cha (parent context expansion) là kỹ thuật bổ sung
nội dung chunk cha cho kết quả chunk con sau khi reranking, nhằm cung cấp thêm
ngữ cảnh cho mô hình sinh mà không làm ảnh hưởng đến quá trình xếp hạng ứng
viên ban đầu.

Bên cạnh đó, HyDE (\textit{Hypothetical Document Embeddings} \cite{hyde2022} --- Nhúng Tài Liệu
Giả Định) đóng vai trò cơ chế dự phòng khi kết quả sau reranking quá ít hoặc
điểm quá thấp. Cụ thể, mô hình ngôn ngữ sinh ra một đoạn văn giả định phù hợp
với câu hỏi, đoạn văn đó được mã hóa thành embedding vector để truy hồi thêm
ứng viên. HyDE chỉ được kích hoạt khi chất lượng truy hồi không đạt ngưỡng tối
thiểu, nhằm tránh tăng độ trễ không cần thiết cho những câu hỏi đã có kết quả
tốt.

% =========================================================
\section{Xếp hạng lại bằng mô hình cross-encoder}
% =========================================================

\subsection{Bi-encoder và cross-encoder}

Retrieval bằng embedding vector hoạt động theo cơ chế bi-encoder (*mã hóa hai
chiều độc lập*): truy vấn và tài liệu được mã hóa tách rời rồi so sánh bằng
khoảng cách vector. Cách này cho tốc độ nhanh và khả năng mở rộng tốt, nhưng
không đánh giá được sự tương tác sâu giữa từng từ của truy vấn và tài liệu.

Ngược lại, cross-encoder (*mã hóa chéo, đánh giá cặp*) nhận đồng thời cặp
(truy vấn, tài liệu) làm đầu vào và cho điểm liên quan chi tiết hơn. Tuy nhiên,
chi phí tính toán cao hơn đáng kể và chỉ phù hợp với tập ứng viên có kích thước
nhỏ. Sự kết hợp bi-encoder (để lấy ứng viên nhanh) và cross-encoder (để rerank
chính xác) là thiết kế phổ biến trong các hệ thống RAG hiện đại.

\begin{table}[H]
\centering
\begin{tabular}{|l|p{4.2cm}|p{4.2cm}|}
\hline
\textbf{Tiêu chí} & \textbf{Bi-encoder} & \textbf{Cross-encoder} \\
\hline
Đầu vào & Truy vấn và tài liệu mã hóa tách rời & Cặp truy vấn--tài liệu mã hóa
cùng lúc \\
\hline
Mục đích & Lấy ứng viên nhanh & Reranking ứng viên \\
\hline
Khả năng mở rộng & Tốt với chỉ mục vector & Phù hợp với tập ứng viên nhỏ \\
\hline
Độ chi tiết & Phụ thuộc khoảng cách vector & Đánh giá độ liên quan trực tiếp \\
\hline
\end{tabular}
\caption{So sánh bi-encoder và cross-encoder}
\label{tab:biencoder_crossencoder}
\end{table}

\subsection{Mô hình reranker BGE}

Hệ thống sử dụng mô hình cross-encoder BGE (BAAI/bge-reranker-v2-m3) để chấm
điểm từng cặp truy vấn--tài liệu sau giai đoạn truy hồi. Mô hình ghi điểm cho từng ứng viên, áp ngưỡng
trước khi cắt top-$k$ và sử dụng cơ chế đảm bảo số tối thiểu để tránh trả
kết quả rỗng khi tất cả điểm nằm dưới ngưỡng. Đáng chú ý, ngưỡng cho văn bản
thường và chunk bảng được cấu hình riêng, vì bảng có cách biểu diễn tuyến tính
khác với văn bản xuôi và thường cho điểm thấp hơn ngay cả khi nội dung liên
quan.

% =========================================================
\section{Hiểu truy vấn và định tuyến}
% =========================================================

\subsection{Phân loại độ phức tạp theo ba tầng}

Không phải câu hỏi nào cũng đòi hỏi tác nhân tự động hay nhiều lượt truy hồi.
Vì vậy, hệ thống phân loại truy vấn thành ba nhóm theo mức độ phức tạp tăng
dần:

\begin{itemize}
    \item \textbf{Hội thoại thông thường}: lời chào hoặc trao đổi xã giao
    ngắn --- hệ thống trả lời trực tiếp mà không cần truy hồi;
    \item \textbf{RAG đơn giản}: câu hỏi học vụ đơn miền hoặc tra cứu chính
    sách đơn giản, đi qua luồng RAG tuyến tính thông thường;
    \item \textbf{RAG phức tạp}: câu hỏi so sánh, cần nhiều nguồn hoặc nhiều
    bước xử lý --- được chuyển sang tác nhân Lập Kế Hoạch--Thực Thi.
\end{itemize}

Cách phân loại ba tầng này giúp hệ thống tránh dùng pipeline nặng cho những
câu hỏi đơn giản, từ đó kiểm soát được độ trễ và chi phí tính toán.

\subsection{Phân loại miền và viết lại truy vấn}

Sau khi xác định độ phức tạp, hệ thống tiếp tục phân loại miền để chọn đúng
kho tri thức cần tìm kiếm. Bộ phân loại sử dụng embedding vector BGE-M3 làm
đặc trưng và phân loại truy vấn thành một hoặc nhiều miền trong bốn kho tri
thức đã định nghĩa. Khi độ tin cậy thấp hoặc khoảng cách giữa hai miền hẹp,
pipeline có thể gọi thêm LLM để hiệu chỉnh kết quả phân loại.

Song song với đó, thành phần viết lại truy vấn (QueryReflector) biến câu hỏi
thực tế thành truy vấn rõ ràng hơn cho retrieval. Cụ thể, nó loại bỏ nhiễu
hội thoại, viết lại câu hỏi thành dạng độc lập khi phát hiện đại từ hoặc tham
chiếu ngầm, đồng thời trích xuất các thực thể như mã ngành, mã học phần và học
kỳ. Một nguyên tắc quan trọng là hệ thống chỉ hợp nhất thông tin hồ sơ người
dùng khi câu hỏi thực sự phụ thuộc vào thông tin đó, nhằm tránh hiện tượng
context bleed đã đề cập ở phần trước.

% =========================================================
\section{RAG có tác nhân tự động}
% =========================================================

\subsection{Khi nào cần tác nhân}

RAG tuyến tính phù hợp khi một lượt truy hồi và một lượt generation đủ để giải
quyết câu hỏi. Tuy nhiên, một số câu hỏi học vụ đòi hỏi nhiều bước xử lý hơn:
so sánh chuẩn ngoại ngữ giữa hai chương trình, kiểm tra điều kiện tốt nghiệp
theo ngành và khóa cụ thể, hay tổng hợp thông tin từ nhiều văn bản khác nhau.
Với những trường hợp như vậy, tác nhân tự động giúp phân rã câu hỏi, lập kế
hoạch retrieval và tổng hợp kết quả một cách có cấu trúc hơn.

\subsection{Mô hình Lập Kế Hoạch--Thực Thi (Planner--Executor)}

Một cách triển khai agentic RAG phổ biến là mô hình Lập Kế Hoạch--Thực Thi
(Planner--Executor). Trong mô hình này, bước Lập Kế Hoạch phân tích câu hỏi và
tạo ra một kế hoạch gồm một số bước nhỏ (giới hạn ở vài bước), mỗi bước xác định
truy vấn con, kho tri thức đích cùng các gợi ý phạm vi như mã ngành hoặc khóa.
Tiếp theo, bước Thực Thi chạy song song các bước này bằng công cụ retrieval và
tự nới lỏng bộ lọc rồi thử lại khi một bước chưa thu được kết quả. Cuối cùng,
bước Tổng Hợp gộp kết quả của toàn bộ các bước và sinh ra câu trả lời. Luồng xử
lý đi một chiều --- lập kế hoạch một lần, thực thi song song rồi tổng hợp --- giúp dễ
kiểm soát độ trễ và chi phí hơn.

Mô hình này phù hợp với bài toán học vụ vì nhiều câu hỏi không thể được giải
quyết bằng một chunk tài liệu duy nhất. Để so sánh hai chương trình đào tạo,
hệ thống cần truy hồi riêng thông tin của từng chương trình rồi mới tổng hợp;
hay để kiểm tra điều kiện cá nhân, hệ thống có thể cần kết hợp thông tin từ
hồ sơ người dùng, quy định đào tạo và chương trình học.

Tuy vậy, tác nhân chỉ nên được kích hoạt khi câu hỏi thực sự cần lập kế hoạch
nhiều bước. Nếu dùng tác nhân cho mọi câu hỏi, hệ thống sẽ tăng độ trễ và chi
phí mà không thu được lợi ích tương xứng. Vì lý do này, hệ thống kết hợp tác
nhân với cơ chế phân loại theo độ phức tạp như đã trình bày ở phần trước.

% =========================================================
\section{Tìm kiếm web dự phòng và tự đánh giá}
% =========================================================

Kho tri thức nội bộ phù hợp với các tài liệu đã được lập chỉ mục, nhưng một
số câu hỏi về lịch đăng ký hoặc thời hạn mới có thể cần dữ liệu cập nhật hơn.
Tìm kiếm web dự phòng (web fallback) là cơ chế tùy chọn bổ sung nguồn từ
Internet, kích hoạt khi nguồn nội bộ không đủ mạnh hoặc câu trả lời bị đánh giá
là thiếu căn cứ. Đây là một cơ chế có thể bật/tắt qua cấu hình và mặc định được
để ở trạng thái thận trọng, chủ yếu nhắm vào nhóm dữ liệu biến động nhanh như
kế hoạch và thông báo. Với miền học vụ, cơ chế này cần được giới hạn ở các nguồn
chính thức để tránh lấy thông tin từ nguồn chưa được kiểm chứng.

Bên cạnh đó, tự đánh giá (self-evaluation) là kỹ thuật sử dụng chính mô hình
hoặc một mô hình đánh giá riêng để kiểm tra câu trả lời sau khi đã sinh ra.
Các tiêu chí thường bao gồm mức độ liên quan, độ bám nguồn và độ đầy đủ. Kết
quả đánh giá có thể giúp hệ thống quyết định thử lại, chuyển sang web fallback
hoặc thông báo thiếu thông tin. Tuy nhiên, đây cũng là một tầng tùy chọn và mặc
định được tắt để tránh tăng độ trễ cho mỗi truy vấn; quan trọng hơn,
self-evaluation chỉ là một lớp kiểm tra bổ sung, không thể thay thế cho bộ đánh
giá ngoại tuyến và kiểm thử hồi quy.

% =========================================================
\section{Lưu trữ, bộ nhớ đệm và bảo mật}
% =========================================================

Một hệ thống trợ lý RAG cần lưu trữ nhiều loại dữ liệu: thông tin người dùng,
hồ sơ học vụ, phiên hội thoại, nhật ký truy vấn, dấu vết tác nhân, tài liệu
tải lên, chunk trung gian, phản hồi và kết quả đánh giá. Cơ sở dữ liệu dạng
tài liệu như MongoDB phù hợp với các trạng thái ứng dụng có cấu trúc linh hoạt.
Trong khi đó, Redis có thể được dùng cho phiên làm việc, bộ nhớ đệm lịch sử
hội thoại và giới hạn tốc độ truy cập nhằm giảm độ trễ cho các truy vấn lặp
lại.

Điều quan trọng là bộ nhớ đệm trong hệ thống RAG cần được thiết kế cẩn thận.
Cụ thể, hệ thống không nên lưu đệm các câu trả lời thiếu nguồn, dùng web
fallback hoặc bị self-evaluation đánh giá thấp. Thêm vào đó, khi tài liệu mới
được lập chỉ mục, các bản ghi đệm liên quan cần được vô hiệu hóa kịp thời để
tránh trả lời dựa trên nguồn cũ.

Về bảo mật, xác thực người dùng giúp bảo vệ lịch sử hội thoại và hồ sơ cá
nhân. Phân quyền giữa sinh viên và quản trị viên là cần thiết để kiểm soát các
chức năng nhạy cảm như tải tài liệu lên, duyệt chunk hay thay đổi cấu hình hệ
thống. Backend cần lấy định danh từ token (*mã xác thực*) thay vì tin vào thông
tin do client gửi trong yêu cầu, nhằm ngăn chặn rủi ro giả mạo.

% =========================================================
\section{Truy vết và đánh giá}
% =========================================================

Do RAG bao gồm nhiều bước trung gian, hệ thống cần ghi lại đầy đủ thông tin
truy vết để hỗ trợ phân tích lỗi. Nhờ các trường nhật ký ghi lại kho tri thức
được chọn, xác suất định tuyến, truy vấn sau viết lại, bộ lọc áp dụng, kết
quả theo kho, điểm reranking và thời gian từng bước, người phát triển có thể
phân biệt được lỗi thuộc về định tuyến sai, bộ lọc quá chặt, retrieval thiếu
nguồn hay generation không bám ngữ cảnh.

Về đánh giá, các chỉ số retrieval thường dùng bao gồm Recall@K (nguồn đúng có
xuất hiện trong top-$K$), MRR@K (vị trí của nguồn đúng đầu tiên) và nDCG@K
(chất lượng thứ tự xếp hạng có chiết khấu logarit). Hệ thống tách thành hai
lớp đánh giá độc lập: đánh giá theo chính sách hiện hành kiểm tra retrieval
trên tài liệu hiện hành có xét đến quan hệ kế thừa, còn đánh giá theo lịch sử
hội thoại kiểm tra khả năng xử lý câu hỏi tiếp nối và hành vi trong các phiên
trước đó. Sự tách biệt này cần thiết vì chính sách học vụ có thể thay đổi theo
thời gian; một câu trả lời lịch sử không còn đúng ở hiện tại không nhất thiết
phản ánh lỗi của hệ thống retrieval hiện hành.

% =========================================================
\section{Kết luận chương}
% =========================================================

Chương này đã trình bày các nền tảng lý thuyết phục vụ trực tiếp cho hệ thống
trợ lý học vụ. Các khái niệm được tổ chức theo trình tự từ bài toán đến giải
pháp: từ đặc thù dữ liệu học vụ đa miền và yêu cầu câu trả lời có căn cứ, đến
nguyên lý RAG, chunking phân cấp, embedding vector đa mô hình, truy hồi lai với
hợp nhất theo thứ hạng đảo nghịch (RRF) và phương án tuyến tính dùng để đối
chứng, lọc siêu dữ liệu, reranking bằng cross-encoder, phân loại độ phức tạp ba
tầng, agentic RAG theo mô hình
Lập Kế Hoạch--Thực Thi, web fallback có kiểm soát và đánh giá hồi quy. Toàn
bộ nền tảng lý thuyết này sẽ được áp dụng cụ thể vào kiến trúc và triển khai
hệ thống trong các chương tiếp theo.

\end{document}